%%%PAGE 298 rot=0 legibility=good summary=剪贴的电磁感应综合题（水平导轨上ab棒与倾斜导轨上ef棒）及其配图，其余为空白
%%%BLOCK id=298-01 ch=12 sec=双棒模型 kind=problem alt=02,06 cont=none y=0.00-0.30
% 页面上部为剪贴的印刷体题目，只有题干与配图，无解答；其余为空白横线页（背面透字已忽略）。剪贴纸右上部有一条手写斜线（疑为裁剪线）
% 手写标记：方框＝“光滑”“质量均为m，长均为L”“动摩擦因数为μ”“不计所有导轨和ab棒的电阻”；椭圆圈＝“同一水平面内，”；下划线＝“最大静摩擦力与滑动摩擦力大小相等，”“一段距离后停止，”“ef棒始终静止”；第(3)问“取走小立柱1和2，且运动过程中”下方另有手写断续下划线（形状不清）
% 配图中另有手写：R_1=R、R/2、θ 等小标注（已含在图里）
\begin{problem}[剪贴题 8]
8．如图所示，金属导轨 $MNC$ 和 $PQD$，$MN$ 与 $PQ$ 平行且间距为 $L$，所在平面与水平面夹角为 $\alpha$，$N$、$Q$ 连线与 $MN$ 垂直，$M$、$P$ 间接有阻值为 $R$ 的电阻；\fbox{\unclear{光滑}}直导轨 $NC$ 和 $QD$ 在\fbox{同一水平面内，}与 $NQ$ 的夹角都为锐角 $\theta$。均匀金属棒 $ab$ 和 $ef$ \fbox{质量均为 $m$，长均为 $L$}，$ab$ 棒初始位置在水平导轨上与 $NQ$ 重合；$ef$ 棒垂直放在倾斜导轨上，与导轨间的\fbox{动摩擦因数为 $\mu$}（$\mu$ 较小），由导轨上的小立柱 1 和 2 阻挡而静止。空间有方向竖直的匀强磁场（图中未画出）。两金属棒与导轨保持良好接触。\fbox{不计所有导轨和 $ab$ 棒的电阻}，$ef$ 棒的阻值为 $R$，\underline{最大静摩擦力与滑动摩擦力大小相等，}忽略感应电流产生的磁场，重力加速度为 $g$。

(1) 若磁感应强度大小为 $B$，给 $ab$ 棒一个垂直于 $NQ$、水平向右的速度 $v_1$，在水平导轨上沿运动方向滑行\underline{一段距离后停止，}\underline{$ef$ 棒始终静止}，求此过程 $ef$ 棒上产生的热量；

(2) 在 (1) 问过程中，$ab$ 棒滑行距离为 $d$，求通过 $ab$ 棒某横截面的电荷量；

(3) 若 $ab$ 棒以垂直于 $NQ$ 的速度 $v_2$ 在水平导轨上向右匀速运动，并在 $NQ$ 位置时取走小立柱 1 和 2，且运动过程中 $ef$ 棒始终静止。求此状态下最强磁场的磁感应强度及此磁场下 $ab$ 棒运动的最大距离。

%FIG p298_a 0.015 0.175 0.375 0.302
\notefig[0.38]{figures/p298_a.png}
\end{problem}
%%%END
%%%PAGEEND 298
